\section{CAD, BIM, simulação e documentação}

\begin{frame}{Ferramentas: cada uma responde a uma pergunta}
\scriptsize
\begin{tabularx}{\textwidth}{>{\bfseries}p{2.6cm}X}
\toprule
Ferramenta & Uso \\
\midrule
CAD elétrico & Plantas, diagramas, detalhes, listas e documentação 2D. \\
BIM & Compatibilização espacial, parâmetros, quantitativos e coordenação multidisciplinar. \\
QElectroTech & Diagramas didáticos/funcionais de comando e potência. \\
ETAP / EasyPower / similares & Fluxo de carga, curto, seletividade, partida, arco e harmônicas conforme módulos. \\
Planilha / Python & Automação de cálculos, auditoria, geração de tabelas e checagens repetitivas. \\
Analisador de energia & Validar modelo com medições reais em retrofit/operação. \\
\bottomrule
\end{tabularx}
\end{frame}

\begin{frame}{Fluxo de simulação elétrica}
\small
\centering
\begin{tikzpicture}[node distance=0.30cm,every node/.style={font=\scriptsize}]
\tikzset{b/.style={draw=corPrim,fill=corDef!7,rounded corners,minimum width=2.35cm,minimum height=0.65cm,align=center}}
\node[b] (a) {Topologia e fontes};
\node[b,right=of a] (b) {Cabos / trafos};
\node[b,right=of b] (c) {Cargas / modos};
\node[b,below=of c] (d) {Fluxo de carga};
\node[b,left=of d] (e) {Curto-circuito};
\node[b,left=of e] (f) {Seletividade / arco};
\node[b,below=of f] (g) {Partida / harmônicas};
\node[b,right=of g] (h) {Relatórios};
\node[b,right=of h] (i) {Revisão do projeto};
\foreach \u/\v in {a/b,b/c,c/d,d/e,e/f,f/g,g/h,h/i}{\draw[-{Stealth},thick,corPrim] (\u)--(\v);}
\draw[-{Stealth},thick,corAccent] (i.east) to[out=0,in=-80] (a.south east);
\end{tikzpicture}
\vspace{2mm}
\begin{caixaAt}
Software calcula o modelo fornecido. Dado errado + modelo errado = resultado preciso e inútil.
\end{caixaAt}
\end{frame}

\begin{frame}{BIM: parâmetros elétricos úteis}
\small
\begin{columns}[T,onlytextwidth]
\column{0.49\textwidth}
\begin{itemize}
\item Circuito e quadro de origem.
\item Potência e tensão.
\item Número de fases.
\item Fator de potência e demanda.
\item Reserva e status do circuito.
\end{itemize}
\column{0.49\textwidth}
\begin{itemize}
\item Comprimento e rota.
\item Cabo/eletroduto/leito.
\item Dispositivo de proteção.
\item Classificação de carga crítica.
\item Código para operação/manutenção.
\end{itemize}
\end{columns}
\begin{caixaObs}
O valor do BIM aumenta quando os parâmetros alimentam cálculo, documentação e operação -- não apenas um modelo 3D bonito.
\end{caixaObs}
\end{frame}

\begin{frame}[fragile]{Automação de checagens com Python}
\scriptsize
\begin{lstlisting}[style=codPy]
from math import sqrt, acos, sin

def corrente_3f(P, Vll, fp=1.0, eta=1.0):
    return P/(sqrt(3)*Vll*fp*eta)

def queda_3f(I, L_m, R_ohm_km, X_ohm_km, fp, Vll):
    phi = acos(fp)
    dv = sqrt(3)*I*(L_m/1000)*(R_ohm_km*fp + X_ohm_km*sin(phi))
    return 100*dv/Vll

I = corrente_3f(42000, 380, 0.88)
dv = queda_3f(I, 60, 0.75, 0.08, 0.88, 380)
print(I, dv)
\end{lstlisting}
\begin{caixaObs}
O script deve ler dados rastreáveis e gerar alerta, não escolher cabo ``por mágica''. O engenheiro continua responsável pelas premissas e validação.
\end{caixaObs}
\end{frame}

\begin{frame}{Documentação para manutenção}
\small
\begin{columns}[T,onlytextwidth]
\column{0.49\textwidth}
\begin{itemize}
\item Unifilar atualizado.
\item Ajustes de relés/disjuntores.
\item Curto-circuito disponível por barra.
\item Identificação de fontes alternativas.
\item Procedimentos de manobra.
\end{itemize}
\column{0.49\textwidth}
\begin{itemize}
\item Etiquetas e riscos aplicáveis.
\item Histórico de intervenções.
\item Ensaios de DR, isolamento e aterramento.
\item Termografias e inspeções.
\item Lista de sobressalentes críticos.
\end{itemize}
\end{columns}
\end{frame}

%=====================================================================
